\documentclass{article}
\usepackage{notes}

\setcounter{section}{0}
\renewcommand{\thesection}{7.\arabic{section}}

\begin{document}
\stdtitle{ANALYSIS II}{Exercise Sheet 7}{Jakob Schildhauer}

\section{Diffeomorphism}
\begin{enumerate}
    \item $(x, y) \mapsto \left(\frac{\arctan x}{\pi} + \frac12. \frac{\arctan y}{\pi} + \frac12\right)$
    \item $f(x)$ is clearly bijective as $f^{-1}(y) =\sqrt[5]{y}$ is defined on the whole of $\R$. However, the derivative of the inverse is not defined at $x = 0$, so $f$ is not a diffeomorphism.

\end{enumerate}

\section{Inverse function I}
$\det J_F(x, y) = 3x^2 y^2 \neq 0 \fa x, y \neq 0$, \\
$(J_F(2, 1)^{-1}) = \frac{1}{12}\begin{pmatrix}
    4 & -4 \\ -1 & 4
\end{pmatrix}$.

\section{Implicit function I}
\begin{enumerate}
    \item \dots 
    \item \dots 
    \item \dots
\end{enumerate}

\section{Multiple choice}
Correct: \titlefont{(a)}, \titlefont{(d)}.

\section{Implicit function II}
\begin{enumerate}
    \item \dots 
    \item \dots 
    \item \dots
\end{enumerate}

\section{Spherical coordinates}
\begin{enumerate}
    \item The images are orthogonal (like a local cartesian coordinate system).
    \item $\Image \Phi = \R^3 \setminus \SET{(x, y, z)}{x \leq 0, y = 0}$
    \item Extensive computation of the determinant where a lot cancels thanks to $\cos^2 + \sin^2$.
    \item Every cartesian point in the image of $\Phi$ has a unique spherical and vice versa, so bijectivity follows. Since $\Phi$ is a local diffeomorphism ($r^2 \sin \theta > 0$), $\Phi$ is a global diffeomorphism. 
\end{enumerate}

\section{Parametrized surfaces}
\begin{enumerate}
    \item Möbius: $n(u, v) = $
    \item 
    \item $(\theta, \varphi) \mapsto (a \sin \theta \cos \varphi, b \sin \theta \sin \varphi, c \cos \theta)$, where $\theta \in (0, \pi), \varphi \in (-\pi, \pi)$ (inheriting the problems of spherical coordinates). 
\end{enumerate}

\section{Graphs as parametrized surfaces}
\begin{enumerate}
    \item $\phi(x, y) = (x, y, f(x, y))$.
    \item 
    \item 
    \item 
\end{enumerate}

\section{The IFT is only a sufficient condition}
\begin{enumerate}
    \item $\frac{\p f(0, 0)}{\p x} = \frac{\p f(0, 0)}{\p y} = 0$
    \item For a fixed $y$, $f$ is strictly decreasing, so there exists a unique solution everywhere. 
    \item $\alpha$  
\end{enumerate}
\end{document}